\begin{claim}{C.83}{Proved}{An actual equivariant infinite-evaluation path without a closed-disk lift}
There is an actual closed disk and a closed,
discrete, free proper interior operator group
$Q\cong\mathbb Z$ satisfying every hypothesis of
C.81, including an actual transitive
boundary-fixed subgroup of its centralizer,
with the following property.

For a specified ordered, locally finite,
$Q$-invariant vertex set $V$, let
$\mathscr E_V$ be the actual image of
simultaneous labelled evaluation
\[
 \operatorname{ev}_V:\mathscr C_Q\longrightarrow\mathscr E_V.
                                                               \tag{C.83.1}
\]
Give $\mathscr E_V$ the subspace topology of
ordered coordinate convergence together with
vague convergence of counting measures in the
specified interior plane coordinate.
The map (C.83.1) is continuous and onto.
There is an explicit based null-homotopic loop in
$\mathscr E_V$ which has no continuous lift
to $\mathscr C_Q$ starting at $1$.
Every point of that loop is the evaluation
of a displayed actual boundary-fixed,
$Q$-commuting homeomorphism. All original
$Q$-labels are retained.

The vertex set can be taken to be the full
lift of a quotient triangulation having the
null-star property of C.82. Thus, for these
exact source and target topologies, infinite
evaluation has neither the path-lifting
property nor the Serre or Hurewicz fibration
property. The proof does not assert a
nontrivial homotopy class in the actual
centralizer.
\end{claim}

\begin{proof}
\emph{The actual disk, operators and transitive source.}
Let $\overline D$ be the closed Euclidean
unit disk, let $O=\operatorname{int}\overline D$,
and use its original Euclidean distance $d$.
The specified interior coordinate and inverse are
\[
 \begin{gathered}
 j:\mathbb R^2\longrightarrow O,\qquad
       j(z)=\frac{z}{1+\|z\|},\\
 j^{-1}:O\longrightarrow\mathbb R^2,\qquad
       j^{-1}(w)=\frac{w}{1-\|w\|}.
 \end{gathered}                                                  \tag{C.83.2}
\]
For $a\in\mathbb R^2$ write $T_a(z)=z+a$.
Set $q=jT_{(1,0)}j^{-1}$ and $Q=\langle q\rangle$,
with exactly these powers as its labels.
The conjugacy (C.83.2) proves freeness and
properness: two compact subsets of the plane
meet under only finitely many integer horizontal
translations. It also proves that $Q$ is a
closed discrete subgroup of $\operatorname{Homeo}(O)$.
Indeed, evaluation at $j(0)$ of any convergent
sequence of its elements bounds the integer
translations in the plane, and a convergent
subsequence of integers is eventually constant.

Each $jT_a j^{-1}$ extends by the identity
on the original circle. We will use the
following explicit estimate repeatedly.
For $z,v\in\mathbb R^2$,
\[
 \|j(z+v)-j(z)\|
       \leq\frac{2\|v\|}{1+\|z\|}.
                                                               \tag{C.83.3}
\]
To prove it, put $r=\|z\|$ and $r'=\|z+v\|$
and write the difference as
$v/(1+r)+(z+v)((1+r')^{-1}-(1+r)^{-1})$.
The reverse triangle inequality bounds the
second summand by $\|v\|/(1+r)$.
The same estimate applies to inverse
translations. Uniform bounds for $a$ on
compact parameter sets therefore give joint
continuity of the extended maps and inverses.

Let $G$ consist of these extended translations.
Their displayed composition law makes $G$
an actual subgroup of $\mathscr C_Q$,
and $G$ is transitive on $O$.
Here $\mathscr C_Q$ is exactly the group in
C.81.1, with metric
\[
 D(h,k)=
 \sup_{w\in\overline D}\|h(w)-k(w)\|
 +\sup_{w\in\overline D}\|h^{-1}(w)-k^{-1}(w)\|.
                                                               \tag{C.83.4}
\]
This is its original compact-open topology
on maps and inverses. In particular, no
interior-only topology is substituted.

\emph{The complete label set and target topology.}
In the plane put
$v_{k,l}=(k/4,l)$ for $(k,l)\in\mathbb Z^2$,
and let $V=\{j(v_{k,l})\}$.
Enumerate $\mathbb Z^2$ by increasing
$\max\{|k|,|l|\}$, using lexicographic order
within each finite shell. This fixes one
literal enumeration by the positive integers.
The operator $q$ acts on its labels by
$(k,l)\mapsto(k+4,l)$.

An ordered configuration is a sequence of
distinct points $z_i\in\mathbb R^2$ for which
every compact planar set contains only
finitely many $z_i$. Write
$\mu_z=\sum_i\delta_{z_i}$ and use
\[
 \begin{split}
 d_{\rm prod}(z,w)
    &=\sum_{i\geq1}2^{-i}\min\{1,\|z_i-w_i\|\},\\
 d_\Sigma(z,w)
    &=d_{\rm prod}(z,w)+d_V(\mu_z,\mu_w),
 \end{split}                                                    \tag{C.83.5}
\]
where $d_V$ metrizes vague convergence:
$\mu_m\to\mu$ means
$\int\varphi\,d\mu_m\to\int\varphi\,d\mu$
for every continuous compactly supported
planar function $\varphi$.
This is the sum topology specified in
\cite[Definition 2.6 and equation (2), pp.~5--6]
      {C_TehConfigurations2026}.
Only that definition is used here.

The actual evaluation is
\[
 \begin{gathered}
 \bigl(\operatorname{ev}_V(h)\bigr)_{k,l}
       =j^{-1}h j(v_{k,l}),\qquad
 \mathscr E_V=\operatorname{ev}_V(\mathscr C_Q),\\
 z_{k+4,l}=z_{k,l}+(1,0)
                \quad(z\in\mathscr E_V).
 \end{gathered}                                                  \tag{C.83.6}
\]
The last equation follows from the original
commutation relation with $q$, not from
discarding its labels. A homeomorphism of
the plane sends this grid to a locally
finite set, because its inverse sends
compact sets to compact sets. Thus every
evaluation in (C.83.6) belongs to the stated
configuration space.

We verify continuity of this exact map.
Suppose $h_m\to h$ in (C.83.4).
Each labelled coordinate converges.
For a compact planar set $K$, uniform
convergence of the inverse maps implies that
all $h_m^{-1}(j(K))$, for sufficiently large
$m$, lie in one compact subset of $O$:
take a small closed neighborhood of the
compact set $h^{-1}(j(K))$ still inside $O$.
Only finitely many vertices of $V$ lie
in that neighborhood. Hence the evaluation
of any fixed compactly supported test
function is eventually a fixed finite
sum of continuous coordinate evaluations.
This proves vague convergence and then
continuity for (C.83.5).

The grid also has the required triangulation
interpretation. Divide its rectangles of
width $1/4$ and height $1$ by the same
diagonal, and transport this plane
triangulation through $j$.
Horizontal translation by $1$ preserves it.
Each closed vertex star has horizontal
width $1/2<1$, so its projection to
$O/Q$ is injective. The planar stars have
one common diameter bound. Formula
(C.83.3) shows that their images have
diameters tending to zero as their
vertices escape in the plane; only
finitely many vertices lie in any fixed
compact set. This is precisely the
finite-large-star property used in C.82.

\emph{The actual up-and-back paths.}
Define the continuous tent function
$b(u)=\max\{0,1-4|u|\}$.
For $n\geq1$ and $s\in[0,1]$ put
\[
 \begin{split}
 F_{n,s}(x,y)&=(x+sn\,b(y-n),y),\\
 F_{n,s}^{-1}(x,y)&=(x-sn\,b(y-n),y).
 \end{split}                                                    \tag{C.83.7}
\]
These formulas prove both inverse identities
and commutation with every $T_{(m,0)}$.
Their displacement is at most $n$.
Consequently (C.83.3) extends
$H_{n,s}=jF_{n,s}j^{-1}$ by the identity
to the original circle, with a jointly
continuous inverse. For each fixed $n$,
$s\mapsto H_{n,s}$ is continuous in
(C.83.4). To check the uniform assertion,
use (C.83.3) outside one large planar disk
and ordinary uniform continuity on its
compact closure. The identical argument
works for the displayed inverse.
Thus $H_{n,s}\in\mathscr C_Q$ and
$H_{n,0}=1$.

Only row $l=n$ of the vertex set moves:
\[
 \bigl(\operatorname{ev}_V(H_{n,s})\bigr)_{k,l}
  =
  \begin{cases}
     (k/4+sn,n),&l=n,\\
     (k/4,l),&l\ne n.
  \end{cases}                                                    \tag{C.83.8}
\]
These equations retain every label in
(C.83.6), including labels in the same
operator orbit.

Let $I_n=[2^{-n-1},2^{-n}]$ and write
\[
 \begin{gathered}
 \lambda_n(t)=2^{n+1}t-1,\qquad
 s_n(t)=2\min\{\lambda_n(t),1-\lambda_n(t)\},
       \quad t\in I_n,\\
 t_n=3\cdot2^{-n-2}.
 \end{gathered}                                                  \tag{C.83.9}
\]
The amplitude is zero at both endpoints
of $I_n$ and one at $t_n$.
Define
\[
 c(t)=
 \begin{cases}
  \operatorname{ev}_V(H_{n,s_n(t)}),&t\in I_n,\ n\geq1,\\
  \operatorname{ev}_V(1),&t=0\text{ or }1/2\leq t\leq1.
 \end{cases}                                                    \tag{C.83.10}
\]
The formulas agree at all overlapping
interval endpoints.

We check both components of the topology
at $t=0$. Any finite collection of fixed
labels avoids the moving row once $n$
is sufficiently large. Its coordinates
then equal their original values for
every $s\in[0,1]$; the remaining product
metric terms are bounded by the tail of
$\sum_i2^{-i}$.
For a compact planar set $K$, its
$y$-coordinates are bounded. For all
sufficiently large $n$, the moving row
and all its images are disjoint from
$K$, since (C.83.7) preserves $y$.
Every test function supported in $K$
therefore has exactly its original
counting-measure evaluation, uniformly
in $s$. This proves vague convergence.
Equivalently, for a metric defined by
a weighted countable test family, first
fix finitely many supports and then
bound the remaining summands by their
weight tail. It proves continuity in
$d_\Sigma$ at zero. Continuity on each
$I_n$ and at its other endpoints has
already been proved. Thus (C.83.10) is
a continuous based loop in the actual
subspace $\mathscr E_V$.
Its based null-homotopy in this same
subspace is the exact map
\[
 C(r,t)=
 \begin{cases}
  \operatorname{ev}_V(H_{n,(1-r)s_n(t)}),
                           &t\in I_n,\ n\geq1,\\
  \operatorname{ev}_V(1),&t=0\text{ or }1/2\leq t\leq1,
 \end{cases}
 \qquad 0\leq r\leq1.
                                                               \tag{C.83.11}
\]
The preceding finite-coordinate and
test-support arguments are uniform in
$r$, since every displayed amplitude
still belongs to $[0,1]$. Together with
continuity on each parameter rectangle
and equality at its interval endpoints,
they prove joint continuity of $C$.
Its faces are $C(0,t)=c(t)$,
$C(1,t)=\operatorname{ev}_V(1)$ and
$C(r,0)=C(r,1)=\operatorname{ev}_V(1)$.

\emph{The lower bound for every proposed lift.}
Suppose a continuous map
$\widetilde c:[0,1]\to\mathscr C_Q$
satisfied
$\widetilde c(0)=1$ and
$\operatorname{ev}_V(\widetilde c(t))=c(t)$.
At the peak time $t_n\to0$, the literal
label $(0,n)$ forces
\[
 \begin{gathered}
 \widetilde c(t_n)j(0,n)=j(n,n),\\
 \sup_{w\in\overline D}
       \|\widetilde c(t_n)w-w\|
 \ \geq\ \|j(n,n)-j(0,n)\|
 \ \geq\ \frac{n}{1+\sqrt2\,n}
 \ >\ \frac13 .
 \end{gathered}                                                  \tag{C.83.12}
\]
The penultimate inequality uses the first
coordinate alone. The last follows from
$\sqrt2<2$ and $n/(1+2n)\geq1/3$.
It holds for every homeomorphism having
the required labelled evaluation,
regardless of its behavior between
vertices. It contradicts continuity at
zero in (C.83.4).

This proves the claimed failure of
path lifting. The lifting axiom of
either named fibration, applied to a
one-point initial space, would supply
precisely that lift, so it also fails.
No contraction of an original compact
sphere family has been disproved:
the loop just constructed lives in
the explicitly stated evaluation-image
topology and has the displayed based
null-homotopy there. In particular this calculation
does not decide the remaining general
centralizer homotopy question in C.12.
\end{proof}
